
\documentclass{article}
\usepackage{colortbl}
\usepackage{makecell}
\usepackage{multirow}
\usepackage{supertabular}

\begin{document}

\newcounter{utterance}

\centering \large Interaction Transcript for game `dond', experiment `coop', episode 16 with Qwen3.6{-}35B{-}A3B{-}FP8{-}without{-}reasoning.
\vspace{24pt}

{ \footnotesize  \setcounter{utterance}{1}
\setlength{\tabcolsep}{0pt}
\begin{supertabular}{c@{$\;$}|p{.15\linewidth}@{}p{.15\linewidth}p{.15\linewidth}p{.15\linewidth}p{.15\linewidth}p{.15\linewidth}}
    \# & $\;$A & \multicolumn{4}{c}{Game Master} & $\;\:$B\\
    \hline

    \theutterance \stepcounter{utterance}  
    & & \multicolumn{4}{p{0.6\linewidth}}{
        \cellcolor[rgb]{0.9,0.9,0.9}{
            \makecell[{{p{\linewidth}}}]{
                \texttt{\tiny{[P1$\langle$GM]}}
                \texttt{Jucați o jocul negocierii cooperative în care trebuie să vă puneți de acord cu un alt jucător cu privire la modul în care se împarte un set de articole între voi.} \\
\\ 
\texttt{Reguli:} \\
\texttt{(a) Tu și celălalt jucător primiți un set de articole. Fiecare dintre voi primește și o funcție de valoare secretă, reprezentând cât apreciezi fiecare tip de obiect.} \\
\texttt{(b) Schimbați mesaje cu celălalt jucător pentru a fi de acord cu privire la cine primește care articole. Fiecare poate trimite maximum 5 mesaje sau termina jocul mai devreme dând o propunere secretă în orice moment.} \\
\texttt{(c) Vi se cere să trimiteți o propunere secretă indicând articolele pe care le doriți, formatate între paranteze pătrate după cum urmează: "{[}Propunere: <număr> <nume articol>, <număr> <nume articol>, <...>{]}"} \\
\texttt{(d) Dacă propunerile tale sunt complementare, adică sunt suficiente articole pentru a îndeplini ambele propuneri, fiecare jucător primește un scor pe baza sumei valorilor articolelor pe care le{-}a primit. În caz contrar, amândoi primiți zero puncte.} \\
\texttt{(e) Scopul tău este să maximizezi suma scorului tău și a scorului celuilalt jucător.} \\
\\ 
\texttt{Să începem.} \\
\\ 
\texttt{Setul de articole disponibile este:} \\
\\ 
\texttt{2 banane, 2 oale, 1 piatră} \\
\\ 
\texttt{Funcția ta de valoare secretă este:} \\
\\ 
\texttt{banană: 4, oală: 1, piatră: 0} \\
\\ 
\texttt{IMPORTANT: Mesajele tale, cu excepția unei propuneri secrete, sunt transmise direct celuilalt jucător, deci nu include niciun răspuns la reguli sau text care anunță mesajul tău. Pentru a face o propunere secretă, utilizează formatul indicat. Nu utiliza paranteze pătrate când comunici cu celălalt jucător, altfel va fi interpretat ca propunerea ta secretă.} \\
            }
        }
    }
    & \\ \\

    \theutterance \stepcounter{utterance}  
    & \multicolumn{4}{p{0.6\linewidth}}{
        \cellcolor[rgb]{0.9,0.9,0.9}{
            \makecell[{{p{\linewidth}}}]{
                \texttt{\tiny{[P1$\rangle$GM]}}
                \texttt{Salut! Sunt de acord să jucăm acest joc.} \\
\\ 
\texttt{Având în vedere că scopul nostru comun este să maximizăm scorul total, cred că o împărțire echilibrată ar fi ideală. Eu apreciez bananele mult mai mult decât oalele, iar piatra nu are valoare pentru mine.} \\
\\ 
\texttt{Propun ca tu să iei ambele oale (care au valoare pentru tine, sper) și piatra, iar eu să iau ambele banane. Astfel, fiecare obiect ajunge la persoana care îl apreciază cel mai mult sau cel puțin nu se risipește.} \\
\\ 
\texttt{Ce părere ai?} \\
            }
        }
    }
    & & \\ \\

    \theutterance \stepcounter{utterance}  
    & & & \multicolumn{2}{p{0.3\linewidth}}{
        \cellcolor[rgb]{0.9,0.9,0.9}{
            \makecell[{{p{\linewidth}}}]{
                \texttt{\tiny{[GM$|$GM]}}
                \texttt{continue} \\
            }
        }
    }
    & & \\ \\

    \theutterance \stepcounter{utterance}  
    & & & \multicolumn{2}{p{0.3\linewidth}}{
        \cellcolor[rgb]{0.9,0.9,0.9}{
            \makecell[{{p{\linewidth}}}]{
                \texttt{\tiny{[GM$|$GM]}}
                \texttt{Salut! Sunt de acord să jucăm acest joc.} \\
\\ 
\texttt{Având în vedere că scopul nostru comun este să maximizăm scorul total, cred că o împărțire echilibrată ar fi ideală. Eu apreciez bananele mult mai mult decât oalele, iar piatra nu are valoare pentru mine.} \\
\\ 
\texttt{Propun ca tu să iei ambele oale (care au valoare pentru tine, sper) și piatra, iar eu să iau ambele banane. Astfel, fiecare obiect ajunge la persoana care îl apreciază cel mai mult sau cel puțin nu se risipește.} \\
\\ 
\texttt{Ce părere ai?} \\
            }
        }
    }
    & & \\ \\

    \theutterance \stepcounter{utterance}  
    & & \multicolumn{4}{p{0.6\linewidth}}{
        \cellcolor[rgb]{0.9,0.9,0.9}{
            \makecell[{{p{\linewidth}}}]{
                \texttt{\tiny{[GM$\rangle$P2]}}
                \texttt{Jucați o jocul negocierii cooperative în care trebuie să vă puneți de acord cu un alt jucător cu privire la modul în care se împarte un set de articole între voi.} \\
\\ 
\texttt{Reguli:} \\
\texttt{(a) Tu și celălalt jucător primiți un set de articole. Fiecare dintre voi primește și o funcție de valoare secretă, reprezentând cât apreciezi fiecare tip de obiect.} \\
\texttt{(b) Schimbați mesaje cu celălalt jucător pentru a fi de acord cu privire la cine primește care articole. Fiecare poate trimite maximum 5 mesaje sau termina jocul mai devreme dând o propunere secretă în orice moment.} \\
\texttt{(c) Vi se cere să trimiteți o propunere secretă indicând articolele pe care le doriți, formatate între paranteze pătrate după cum urmează: "{[}Propunere: <număr> <nume articol>, <număr> <nume articol>, <...>{]}"} \\
\texttt{(d) Dacă propunerile tale sunt complementare, adică sunt suficiente articole pentru a îndeplini ambele propuneri, fiecare jucător primește un scor pe baza sumei valorilor articolelor pe care le{-}a primit. În caz contrar, amândoi primiți zero puncte.} \\
\texttt{(e) Scopul tău este să maximizezi suma scorului tău și a scorului celuilalt jucător.} \\
\\ 
\texttt{Să începem.} \\
\\ 
\texttt{Setul de articole disponibile este:} \\
\\ 
\texttt{2 banane, 2 oale, 1 piatră} \\
\\ 
\texttt{Funcția ta de valoare secretă este:} \\
\\ 
\texttt{banană: 2, oală: 0, piatră: 6} \\
\\ 
\texttt{IMPORTANT: Mesajele tale, cu excepția unei propuneri secrete, sunt transmise direct celuilalt jucător, deci nu include niciun răspuns la reguli sau text care anunță mesajul tău. Pentru a face o propunere secretă, utilizează formatul indicat. Nu utiliza paranteze pătrate când comunici cu celălalt jucător, altfel va fi interpretat ca propunerea ta secretă.} \\
\\ 
\texttt{Salut! Sunt de acord să jucăm acest joc.} \\
\\ 
\texttt{Având în vedere că scopul nostru comun este să maximizăm scorul total, cred că o împărțire echilibrată ar fi ideală. Eu apreciez bananele mult mai mult decât oalele, iar piatra nu are valoare pentru mine.} \\
\\ 
\texttt{Propun ca tu să iei ambele oale (care au valoare pentru tine, sper) și piatra, iar eu să iau ambele banane. Astfel, fiecare obiect ajunge la persoana care îl apreciază cel mai mult sau cel puțin nu se risipește.} \\
\\ 
\texttt{Ce părere ai?} \\
            }
        }
    }
    & \\ \\

    \theutterance \stepcounter{utterance}  
    & & & \multicolumn{4}{p{0.6\linewidth}}{
        \cellcolor[rgb]{0.9,0.9,0.9}{
            \makecell[{{p{\linewidth}}}]{
                \texttt{\tiny{[GM$\langle$P2]}}
                \texttt{Salut! Mulțumesc pentru propunerea inițială.} \\
\\ 
\texttt{Trebuie să clarificăm o mică neînțelegere: în descrierea jocului, am menționat că funcția mea de valoare este banană: 2, oală: 0, piatră: 6. Deoarece scopul nostru este să maximizăm scorul total (suma scorurilor noastre), ar fi ideal ca fiecare să primească articolele pe care le valorizăm cel mai mult.} \\
\\ 
\texttt{Eu apreciez piatra foarte mult (6 puncte), în timp ce bananele au o valoare moderată pentru mine (2 puncte), iar oalele nu au valoare (0 puncte).} \\
\\ 
\texttt{Propunerea ta inițială îmi dă mie oalele și piatra. Deoarece oalele nu au valoare pentru mine, acest lucru înseamnă că eu aș primi doar piatra (6 puncte), iar tu ai primi ambele banane (4 puncte totale pentru tine, dacă valoarea ta pentru banană e similară sau mai mare). Totuși, dacă eu aș primi și oalele, nu câștig nimic suplimentar, dar ocup loc în "portofelul" meu de articole fără beneficiu.} \\
\\ 
\texttt{Mai important, dacă eu primesc piatra (6) și tu primești bananele (să zicem 2 fiecare = 4), scorul total este 10. Dar dacă eu primesc piatra și oalele (6+0=6) și tu primești bananele (4), scorul total rămâne 10. Nu există un câștig suplimentar pentru mine din a lua oalele.} \\
\\ 
\texttt{Totuși, pentru a fi corecți și a maximiza eficiența, propun următoarea împărțire:} \\
\texttt{{-} Eu primesc: 1 piatră} \\
\texttt{{-} Tu primești: 2 banane, 2 oale} \\
\\ 
\texttt{De ce? Pentru că eu valorizez piatra la 6, iar tu, presupunând că apreciezi bananele (și poate oalele, chiar dacă nu știu valoarea ta), vei primi articolele care ți se potrivesc. Oalele au valoare 0 pentru mine, deci nu îmi pasă dacă le iei tu. Bananele au valoare 2 pentru mine, dar dacă le iei tu, eu pierd 2 puncte potențiale, dar tu câștigi. Dacă eu aș lua o banană, aș avea 6+2=8, iar tu ai 2+0=2 (dacă oalele au valoare 0 și pentru tine). Dar dacă tu iei bananele, tu ai 4 (sau mai mult), iar eu am 6. Total 10.} \\
\\ 
\texttt{Aștept propunerea ta secretă sau un contra{-}argument. Dacă ești de acord cu împărțirea: eu iau piatra, tu iei bananele și oalele, te rog să îmi trimiți o propunere secretă care să reflecte acest lucru pentru tine (adică să ceri bananele și oalele).} \\
\\ 
\texttt{Până atunci, iată propunerea mea secretă pentru a începe procesul de aliniere:} \\
\\ 
\texttt{{[}Propunere: 1 piatră{]}} \\
            }
        }
    }
     \\ \\

    \theutterance \stepcounter{utterance}  
    & & & \multicolumn{2}{p{0.3\linewidth}}{
        \cellcolor[rgb]{0.9,0.9,0.9}{
            \makecell[{{p{\linewidth}}}]{
                \texttt{\tiny{[GM$|$GM]}}
                \texttt{proposal} \\
            }
        }
    }
    & & \\ \\

    \theutterance \stepcounter{utterance}  
    & & & \multicolumn{2}{p{0.3\linewidth}}{
        \cellcolor[rgb]{0.9,0.9,0.9}{
            \makecell[{{p{\linewidth}}}]{
                \texttt{\tiny{[GM$|$GM]}}
                [0, 0, 1]
            }
        }
    }
    & & \\ \\

    \theutterance \stepcounter{utterance}  
    & & \multicolumn{4}{p{0.6\linewidth}}{
        \cellcolor[rgb]{0.9,0.9,0.9}{
            \makecell[{{p{\linewidth}}}]{
                \texttt{\tiny{[P1$\langle$GM]}}
                \texttt{Celălalt jucător a trimis o propunere secretă. Acum este rândul tău să trimiți o propunere proprie.} \\
            }
        }
    }
    & \\ \\

    \theutterance \stepcounter{utterance}  
    & \multicolumn{4}{p{0.6\linewidth}}{
        \cellcolor[rgb]{0.9,0.9,0.9}{
            \makecell[{{p{\linewidth}}}]{
                \texttt{\tiny{[P1$\rangle$GM]}}
                \texttt{{[}Propunere: 2 banane{]}} \\
            }
        }
    }
    & & \\ \\

    \theutterance \stepcounter{utterance}  
    & & & \multicolumn{2}{p{0.3\linewidth}}{
        \cellcolor[rgb]{0.9,0.9,0.9}{
            \makecell[{{p{\linewidth}}}]{
                \texttt{\tiny{[GM$|$GM]}}
                \texttt{proposal} \\
            }
        }
    }
    & & \\ \\

    \theutterance \stepcounter{utterance}  
    & & & \multicolumn{2}{p{0.3\linewidth}}{
        \cellcolor[rgb]{0.9,0.9,0.9}{
            \makecell[{{p{\linewidth}}}]{
                \texttt{\tiny{[GM$|$GM]}}
                [2, 0, 0]
            }
        }
    }
    & & \\ \\

    \theutterance \stepcounter{utterance}  
    & & & \multicolumn{2}{p{0.3\linewidth}}{
        \cellcolor[rgb]{0.9,0.9,0.9}{
            \makecell[{{p{\linewidth}}}]{
                \texttt{\tiny{[GM$|$GM]}}
                [[2, 0, 0], [0, 0, 1]]
            }
        }
    }
    & & \\ \\

\end{supertabular}
}

\end{document}
